\scanpage{123}\mspage{68}
La première relation se vérifie tautologiquement, elle revient à dire que dans $\unsure{A|\mathcal A}\Sigma^{-1}$, on a un composé identité de
\[
 X\overset{i_X}{\unsure{\text{---}|\longrightarrow}}C(\unsure{x|X})\overset{i_X^{-1}}{\text{---}}X\qquad !
\]
La seconde signifie que pour \unsure{tt|$\forall$} $x$, on a un composé identité dans $A\Sigma_0^{-1}$
\[
 C(X)\xrightarrow{i_{C(X)}}C(C(X))\xrightarrow{C(p_X)}C(X)
\]
Or $p_X:C(X)\text{---}X$ avait été défini dans $A\Sigma^{-1}$ \deleted{\unsure{que}} comme inverse $i_X^{-1}$ de $i_X$, d'autre part le \deleted{$\unsure{C}$} \inserted{foncteur $C$ sur $A$ \unsure{définit|défini} par \unsure{hyp}} $A\Sigma^{-1}\text{---}A\Sigma_0^{-1}$, plus précisément $C(\Sigma)\subset\Sigma_0$, ainsi $C(i_X)$ est inversible dans \deleted{$\unsure{C\,C}$} $A\Sigma_0^{-1}$, et $C(p_X)=C(i_X)^{-1}$. Donc la compatibilité revient à la relation
\[
 C(i_X)^{-1}\deleted{\unsure{C}}\,i_{C(X)}=\mathrm{id}\qquad\text{\uline{dans $A\Sigma_0^{-1}$}}
\]
i.e.
\[
 \boxed{i_{C(X)}=C(i_X)}:\quad C(X)\longrightarrow C^2(X).
\]
En d'autres termes, il faut prouver le

\noindent\begin{tabular}{@{}|p{.94\linewidth}@{}}
\uline{Lemme} \quad Soit $X$ dans $\mathrm{Cat}_{/S}$, considérons les deux flèches
\[
 i_{C(X)},\ C(i_X):\quad C(X)\unsure{\rightrightarrows|\substack{\text{---}\\\text{---}}}C(C(X)).
\]
Ces flèches sont égales, sinon \deleted{\unsure{du}}au sens absolu, du moins dans $(\mathrm{Cat}_{/S}\unsure{]|)}(W_S^d)^{-1}$.
\end{tabular}

\deleted{\unsure{Défi} \unsure{de}}
\[
 \begin{aligned}
 C(X)&=\{x,c\mid x\in\mathrm{Ob}\,X,\ c:f_X(x)\underset{\substack{\text{chemin}\\\text{dans }S}}{\longrightarrow}t\}\\
 &\hspace{4em}(f_X:X\unsure{\longrightarrow|\text{---}} S\ \text{la proj. str.}\\
 C^2(X)&=\{x,c,c'\mid x\in\mathrm{Ob}\,X,\ \deleted{\unsure{c:}}f_X(x)\xrightarrow{c}t\overset{c'}{\longrightarrow}t'\}\\
 &\hspace{4em} c,c'\ \text{chemins dans }S,
 \end{aligned}
\]
les projections dans $S$ sont données par
\[
 \begin{aligned}
 (x,c)&\longmapsto t=\mathrm{but}\,c\\
 (x,c,c')&\longmapsto t'=\mathrm{but}\,c'.
 \end{aligned}
\]
% archive sheet number: 122
